\documentclass[11pt,a4paper]{article}
\usepackage[utf8]{inputenc}
\usepackage[spanish]{babel}
\usepackage{amsmath}
\usepackage{amsfonts}
\usepackage{amssymb}
\usepackage{array}
\usepackage{geometry}
\usepackage{booktabs}
\geometry{
    a4paper,
    total={185mm,260mm},
    left=12mm,
    top=15mm,
}

\setlength{\parindent}{0pt}
\renewcommand{\arraystretch}{2.2}

\begin{document}

\begin{center}
    {\LARGE\bfseries HOJA DE REFERENCIA} \\[3mm]
    {\Large\bfseries ARQUITECTURA DEL COMPUTADOR} \\[2mm]
    {\large\bfseries PRESTACIONES} \\[10mm]

    {\large\bfseries 1. FORMULARIO CRÍTICO DE RENDIMIENTO} \\[5mm]

    \large 
    \begin{tabular}{|m{4.5cm}|m{5.5cm}|m{7.5cm}|}
    \hline
    \centering \textbf{Nombre de la Fórmula} & \centering \textbf{Fórmula} & \centering \textbf{Motivo} \cr \hline
    
    Relación de Rendimiento & 
    \centering $\displaystyle n = \frac{\text{Rend}_X}{\text{Rend}_Y} = \frac{\text{Tiempo}_Y}{\text{Tiempo}_X}$ & 
    Se usa para comparar dos máquinas y saber cuántas veces ($n$) una es más rápida que la otra. \cr \hline
    
    Tiempo de CPU (Frecuencia) & 
    \centering $\displaystyle \text{T}_{\text{CPU}} = \frac{I \times \text{CPI}}{f}$ & 
    Se aplica para calcular el tiempo total en segundos cuando conocemos la frecuencia en Hz, MHz o GHz. \cr \hline
    
    Tiempo de CPU (Ciclo) & 
    \centering $\displaystyle \text{T}_{\text{CPU}} = I \times \text{CPI} \times T_c$ & 
    Se aplica para calcular el tiempo total si conocemos la duración de un ciclo de reloj (ns o ps). \cr \hline
    
    CPI Promedio & 
    \centering $\displaystyle \text{CPI}_{\text{Prom}} = \sum (\text{CPI}_i \times F_i)$ & 
    Se usa cuando el programa tiene distintos tipos de instrucciones con diferentes costes de ciclos. \cr \hline
    
    Ley de Amdahl & 
    \centering $\displaystyle G_{\text{Global}} = \frac{1}{(1 - f_m) + \frac{f_m}{G_m}}$ & 
    Sirve para calcular cuánto mejora el rendimiento total al acelerar solo una parte específica del sistema. \cr \hline
    
    MIPS & 
    \centering $\displaystyle \text{MIPS} = \frac{I}{\text{Tiempo} \times 10^6}$ & 
    Se utiliza para medir el rendimiento en términos de millones de instrucciones ejecutadas por segundo. \cr \hline
    \end{tabular}
\end{center}

\vspace*{\fill}
\begin{flushright}
    \textbf{\large Gabriela Mora} \\
    \textbf{\large V-31.541.893}
\end{flushright}
\newpage

\begin{center}
    {\large\bfseries 2. SISTEMAS DE UNIDADES Y CONVERSIÓN} \\[8mm]
    
    \large 
    
    \textbf{\large A) Frecuencia (Hertz - Hz)} \\[3mm]
    \begin{tabular}{|m{5cm}|m{4cm}|m{7cm}|}
    \hline
    \centering \textbf{Unidad} & \centering \textbf{Símbolo} & \centering \textbf{Multiplicador matemático} \cr \hline
    KiloHertz & \centering KHz & \centering $\times 10^3$ (1.000 Hz) \cr \hline
    MegaHertz & \centering MHz & \centering $\times 10^6$ (1.000.000 Hz) \cr \hline
    GigaHertz & \centering GHz & \centering $\times 10^9$ (1.000.000.000 Hz) \cr \hline
    \end{tabular}
    
    \vspace{12mm}
    
    \textbf{\large B) Tiempo (Segundos - s)} \\[3mm]
    \begin{tabular}{|m{5cm}|m{4cm}|m{7cm}|}
    \hline
    \centering \textbf{Unidad} & \centering \textbf{Símbolo} & \centering \textbf{Equivalencia en segundos} \cr \hline
    Milisegundo & \centering ms & \centering $\times 10^{-3}$ (0,001 s) \cr \hline
    Microsegundo & \centering $\mu$s & \centering $\times 10^{-6}$ (0,000001 s) \cr \hline
    Nanosegundo & \centering ns & \centering $\times 10^{-9}$ (0,000000001 s) \cr \hline
    Picosegundo & \centering ps & \centering $\times 10^{-12}$ (0,000000000001 s) \cr \hline
    \end{tabular}
    
    \vspace{12mm}
    
    \textbf{\large C) Capacidad (Bytes - B)} \\[3mm]
    \begin{tabular}{|m{5cm}|m{4cm}|m{7cm}|}
    \hline
    \centering \textbf{Unidad} & \centering \textbf{Símbolo} & \centering \textbf{Tamaño Exacto (Base 2)} \cr \hline
    KiloByte & \centering KB & \centering $2^{10}$ Bytes (1.024 B) \cr \hline
    MegaByte & \centering MB & \centering $2^{20}$ Bytes (1.048.576 B) \cr \hline
    GigaByte & \centering GB & \centering $2^{30}$ Bytes (1.073.741.824 B) \cr \hline
    \end{tabular}
\end{center}

\vspace*{\fill}
\begin{flushright}
    \textbf{\large Gabriela Mora} \\
    \textbf{\large V-31.541.893}
\end{flushright}

\end{document}